\subsection{The two Dirichlet characters attached to a host}
\label{sec:host-xpp-character-input}

% BEGIN HXD1
\paragraph{HXD1. Splitting the actual host character.}
Fix the host's algebraic closure $\Omega$ of $\mathbb Q$, its odd
prime $\ell$, its coefficient DVR $O=c.O$, its level $N>0$, its
exponent $e$, and its Dirichlet character $\psi$ of level $N$.
Write $k=O/\mathfrak m_O$. The ring $O$ is finite over
$\mathbb Z_\ell$ and contained in $\overline{\mathbb Q}_\ell$.
Its residue field is finite of characteristic $\ell$. Indeed, a
maximal ideal of a finite integral extension contracts to a maximal
ideal, so $\ell O\subseteq\mathfrak m_O$, and $O/\ell O$ is finite
over $\mathbb F_\ell$. In this section a character evaluated on a
unit is identified with its unit value, and the bar denotes reduction
to $k$.

Split $\psi=\theta_0\kappa$ into its prime-to-$\ell$ and
$\ell$-power parts in the finite cyclic group generated by $\psi$.
In the CRT decomposition of $(\mathbb Z/N)^\times$ let $\psi_2$ be
the component on the factor at $2$, inflated to level $N$. That factor
is a $2$-group, hence $\kappa$ is trivial on it. The $2$-component
of $\theta_0$ equals $\psi_2$. Thus the source character
$\theta=\operatorname{th1}(\ell,\psi)$, obtained by dividing
$\theta_0$ by its $2$-component, satisfies
\[
     \psi=\psi_2\theta\kappa,\qquad
     \ell\nmid\operatorname{ord}(\theta),\qquad
     \operatorname{cond}(\theta^2)\text{ is odd}.
\]
The conductor assertion follows because $\theta$ factors through the
odd part of $N$. Reduction sends $\kappa$ to $1$ on all units, since
$X^{\ell^s}-1=(X-1)^{\ell^s}$ in $k[X]$ for any $\ell^s$ killing
$\kappa$. Put
\[
 t=c_2(\psi),\quad n=t-1,\quad
 d=\operatorname{ord}_{\operatorname{lcm}(4,\operatorname{ord}\theta)}(\ell),
 \qquad L=\ell^d-1.
\]
Here $\operatorname{cond}(\psi_2)=2^t$ and $2^t\mid N$.
The multiplicative order exists since $\ell$ is coprime to the
modulus. Therefore $d>0$, $4\mid L$, $\operatorname{ord}(\theta)\mid L$,
and $\ell-1\mid L$ by the geometric sum formula.
The host's $n_0(\ell,\psi)$ equals $v_2(L)$ by definition.
% END HXD1

% BEGIN HXD2
\paragraph{HXD2. The primitive pair and the conductor bound.}
Use the Teichm\"uller character $\omega:(\mathbb Z/\ell)^\times
\to\mathbb Z_\ell^\times\to O^\times$. The unique lift of each
nonzero residue to a root of $X^{\ell-1}-1$ exists by Hensel's
lemma and is multiplicative by uniqueness. Thus $\omega^{\ell-1}=1$,
and its residual Galois character equals the reduction of the
$\ell$-adic cyclotomic character. On the common units take
$\varphi$ to be the primitive character inducing
$\omega^{2e}\theta^2$. If $C=\operatorname{cond}(\theta^2)$,
the product is already defined at level $\ell C$ by inflating
the two factors. Hence its conductor $f$ divides $\ell C$,
whether or not $\ell\mid C$. In particular $f>0$ is odd.
HXD1 shows that both factors are killed by $L$. The induced primitive
character is also killed by $L$, by surjectivity of reduction on unit
groups to each divisor of the modulus. Moreover $\varphi(-1)=1$
since both factors are squares at $-1$.

Let $\psi_2^\circ$ be the unique primitive character of conductor
$2^t$ inducing $\psi_2$. On units modulo $2^t$ take
$(\psi_2^\circ)^2\chi_4$, where $\chi_4$ is the nontrivial
character of $(\mathbb Z/4)^\times$ in $O^\times$. Its values on
units modulo $N$ agree with $\psi_2^2\chi_4$.
The host assumption $t\ge4$ gives
$(\mathbb Z/2^t)^\times=\langle-1\rangle\times\langle5\rangle$,
where $5$ has order $2^{t-2}$. The last assertion follows by
induction from $v_2(5^{2^j}-1)=j+2$, and the powers of $5$ are
precisely the classes $1\bmod4$. The kernel of reduction from
$2^t$ to $2^{t-1}$ is generated by $5^{2^{t-3}}$. Primitivity
at $2^t$ forces $\psi_2^\circ(5)$ to have exact order $2^{t-2}$.
Consequently $((\psi_2^\circ)^2\chi_4)(5)$ has exact order $2^{t-3}$,
as $\chi_4(5)=1$. The product descends to level $2^{t-1}$.
It cannot descend to $2^{t-2}$, where the order of the
class of $5$ is $2^{t-4}$, including $t=4$. This proves primitivity
of the resulting character $\xi$ at $2^n$. At $-1$ it takes value $-1$, giving
$\varphi(-1)\xi(-1)=-1$.

The input witness $q$ of the host supplies precisely $q>0$, $C\mid q$,
and $2(\ell q)4^{n_0(\ell,\psi)}\le 2^t$.
Since $f\mid\ell C\mid\ell q$, cancellation of $2$ gives
\[
       f\,4^{v_2(L)}\le(\ell q)4^{v_2(L)}\le2^{t-1}=2^n .
\]
Together with HXD1 this verifies every field of the Bernoulli
$\operatorname{XData}(\ell,d,\varphi,\xi)$, with the original
$\ell$, the prescribed $e$, and both characters valued in $O$.
% END HXD2

% BEGIN HXD3
\paragraph{HXD3. The identity on the absolute Galois group.}
Let $M=\ell N^2$. Since $C\mid N$ we have $f\mid\ell N$, and
$2^n\mid N$, so $f2^n\mid M$; also $4$, $\ell$, and $N$ divide $M$.
For $g\in G_{\mathbb Q}=\operatorname{Gal}(\Omega/\mathbb Q)$,
choose an integer $a$ representing its action on $\mu_M$. Then
$(a,M)=1$. Restricting the same action to $\mu_b$ for each $b\mid M$
gives $a\bmod b$. Thus evaluation of every character in the computation
below uses the same $a$. On these units the definition of the product
at the common level gives
\[
  (\operatorname{prodX}(\varphi,\xi))(a)=
        \big(\omega(a)^e\psi_2(a)\theta(a)\big)^2\chi_4(a).
\]
Write $\chi_\ell$ for the $O$-valued cyclotomic character and
$\psi_G$ for the Galois character induced by $\psi$.
By HXD1, $\bar\kappa(a)=1$. By HXD2,
$\overline{\omega(a)}=\bar\chi_\ell(g)$. The defining equations
for the supplied host frame give
$\Psi(\pi g)=\chi_\ell(g)^e\psi_G(g)$ and
$\eta(\pi g)=\chi_4(a)$. We obtain the equality of $k^\times$-valued
characters
\[
   \operatorname{redChar}(\operatorname{prodX}(\varphi,\xi))(g)
     =\bar\Psi(\pi g)^2\eta(\pi g).
\]
Here and below $\eta$ on the right is sent from $\{\pm1\}$ to
$k^\times$. In particular BR7 permits the inverse character
$\gamma(g)=\bar\Psi(\pi g)^{-2}\eta(\pi g)$, since
$\eta(\pi g)^{-1}=\eta(\pi g)$. Its coefficient field and reduction
map are $k$ and $O\twoheadrightarrow k$, and $f2^n>0$ as required.
% END HXD3
